% ============================================
% Введение
% ============================================

\chapter*{ВВЕДЕНИЕ}
\addcontentsline{toc}{chapter}{ВВЕДЕНИЕ}

\section*{Актуальность исследования}
\addcontentsline{toc}{section}{Актуальность исследования}

Системная динамика как научная дисциплина и методология моделирования сложных систем имеет глубокие исторические корни, уходящие к работам Джей Форрестера в 1960-х годах в Массачусетском технологическом институте. За прошедшие десятилетия методология зарекомендовала себя как мощный инструмент анализа нелинейного поведения сложных систем во времени, получив широкое распространение в областях управления предприятиями, экологического моделирования, анализа социально-экономических процессов и формирования государственной политики.

Традиционно создание системно-динамических моделей требует от исследователя существенных затрат времени и интеллектуальных ресурсов. Процесс включает в себя этапы качественного концептуального моделирования, построения диаграмм причинно-следственных связей (Causal Loop Diagrams, CLD), разработки диаграмм потоков и уровней (Stock and Flow Diagrams, SFD), формализации уравнений модели, параметризации, верификации и валидации. Каждый из этих этапов требует глубокого понимания как предметной области моделируемой системы, так и специфики самой методологии системной динамики.

Существенным барьером для широкого внедрения системной динамики остаётся необходимость в специализированной подготовке. Исследователям необходимо освоить специфические инструменты моделирования, такие как Vensim, Stella или AnyLogic, понять семантику конструкций системной динамики, а также развить навыки перехода от текстового описания системы к формальной модели. Это ограничивает применение методологии кругом профессиональных системных аналитиков и замедляет распространение системно-динамического подхода в смежных дисциплинах.

Параллельно с развитием методологии системной динамики в последние годы наблюдается бурный прогресс в области больших языковых моделей (Large Language Models, LLM). Модели семейства GPT, Claude, а также открытые альтернативы (LLaMA, Mistral, DeepSeek) продемонстрировали способность к пониманию сложных текстовых описаний, генерации структурированных данных и формальному представлению знаний. Эти способности создают предпосылки для использования LLM в качестве интеллектуального посредника между человеком и инструментами системной динамики.

Актуальность настоящего исследования обусловлена конвергенцией двух указанных трендов: зрелости методологии системной динамики, нуждающейся в инструментах снижения барьера входа, и достижением большими языковыми моделями уровня компетентности, достаточного для участия в процессе моделирования. Создание программной платформы, интегрирующей LLM с конвейером системно-динамического моделирования, представляет собой направление, имеющее как научную новизну, так и значительный практический потенциал.

\section*{Цель и задачи исследования}
\addcontentsline{toc}{section}{Цель и задачи исследования}

\textbf{Целью} настоящей научно-исследовательской работы является разработка методов и алгоритмов автоматизированной генерации системно-динамических моделей на основе текстовых описаний с использованием больших языковых моделей, а также создание программной платформы \texttt{sdCloud}, реализующей предложенные подходы.

Для достижения поставленной цели сформулированы следующие \textbf{задачи}:

\begin{enumerate}
    \item Провести обзор предметной области: изучить теоретические основы системной динамики, проанализировать существующие программные платформы моделирования, исследовать возможности больших языковых моделей в контексте генерации структурированных данных, а также проанализировать существующие работы по интеграции LLM и системной динамики.

    \item Спроектировать архитектуру программной платформы \texttt{sdCloud}, включающей подсистемы взаимодействия с пользователем, управления сессиями моделирования, интеграции с LLM, генерации и валидации моделей, а также хранения данных.

    \item Реализовать backend-подсистему на базе фреймворка FastAPI с использованием реляционной базы данных PostgreSQL, кэширования Redis и асинхронной очереди задач Taskiq.

    \item Реализовать frontend-приложение на базе фреймворка Next.js с применением методологии Feature-Sliced Design (FSD) и библиотеки компонентов shadcn/ui.

    \item Разработать конвейер генерации системно-динамических моделей, включающий этапы извлечения переменных, создания концептуальной карты (CLD), формирования системы уравнений, построения диаграммы потоков и уровней (SFD), а также генерации исполняемого Python-кода и Vensim-экспорта.

    \item Реализовать модуль валидации графовых структур моделей с использованием библиотеки NetworkX для проверки корректности топологии сгенерированных диаграмм и согласованности переменных, уравнений и CLD.

    \item Разработать детерминированный генератор Python-кода по структурированным уравнениям и модуль fallback-симуляции на базе SciPy \texttt{solve\_ivp} для случаев, когда прямая загрузка в PySD невозможна.

    \item Реализовать гостевой режим работы платформы, позволяющий анонимным пользователям генерировать модели и сохранять историю взаимодействия в браузере без обращений к серверной базе данных.

    \item Провести тестирование и оценку работоспособности разработанной платформы на типовых задачах системной динамики.
\end{enumerate}

\section*{Объект и предмет исследования}
\addcontentsline{toc}{section}{Объект и предмет исследования}

\textbf{Объектом исследования} являются процессы создания и модификации системно-динамических моделей с применением методов искусственного интеллекта.

\textbf{Предметом исследования} выступают методы и алгоритмы автоматизированной генерации системно-динамических моделей на основе текстовых описаний с использованием больших языковых моделей, а также архитектурные решения программных платформ, поддерживающих подобную генерацию.

\section*{Научная новизна}
\addcontentsline{toc}{section}{Научная новизна}

Научная новизна работы заключается в следующем:

\begin{itemize}
    \item Предложен многоэтапный конвейер генерации системно-динамических моделей из естественно-языковых описаний, включающий промежуточные представления в виде CLD и SFD, что обеспечивает прозрачность и интерпретируемость процесса генерации.

    \item Разработан метод валидации графовых структур системно-динамических моделей, основанный на анализе топологических свойств графа причинно-следственных связей с использованием алгоритмов на графах, а также на проверке согласованности между переменными, уравнениями и CLD.

    \item Предложен детерминированный метод генерации исполняемого Python-кода по структурированным уравнениям модели, заменяющий нестабильную генерацию кода непосредственно LLM и снижающий влияние галлюцинаций языковой модели на корректность симуляции.

    \item Разработан fallback-механизм симуляции на основе SciPy \texttt{solve\_ivp}, позволяющий выполнять модели в случаях несовместимости сгенерированного кода с текущей версией PySD.

    \item Создана программная платформа \texttt{sdCloud}, интегрирующая современные веб-технологии, базы данных, асинхронную обработку задач, API больших языковых моделей и гостевой режим работы в единое решение для AI-ассистированного системно-динамического моделирования.
\end{itemize}

\section*{Практическая значимость}
\addcontentsline{toc}{section}{Практическая значимость}

Практическая значимость результатов работы определяется следующими факторами:

\begin{itemize}
    \item Разработанная платформа \texttt{sdCloud} может быть использована в образовательном процессе для обучения основам системной динамики, позволяя студентам фокусироваться на анализе моделей, а не на технических аспектах их создания.

    \item Платформа снижает барьер входа в методологию системной динамики для специалистов, не имеющих специализированной подготовки в области моделирования, в том числе благодаря гостевому режиму без обязательной регистрации.

    \item Детерминированный генератор Python-кода и fallback-симуляция через SciPy повышают надёжность и воспроизводимость результатов по сравнению с прямой генерацией исполняемого кода LLM.

    \item Исходный код платформы открыт и может служить основой для дальнейших исследований в области автоматизированного моделирования сложных систем.

    \item Архитектурные решения, применённые в проекте, могут быть адаптированы для создания аналогичных систем в смежных предметных областях.
\end{itemize}

\section*{Структура отчёта}
\addcontentsline{toc}{section}{Структура отчёта}

Отчёт о научно-исследовательской работе состоит из введения, четырёх глав, заключения и списка литературы.

В \textbf{первой главе} представлен обзор литературы по теоретическим основам системной динамики, существующим программным платформам моделирования, большим языковым моделям и существующим работам по интеграции LLM и системной динамики.

Во \textbf{второй главе} описана архитектура программной платформы \texttt{sdCloud}, включая подсистемы backend, frontend, модули генерации и валидации моделей, интеграцию с внешними сервисами LLM, гостевой режим и механизм fallback-симуляции.

В \textbf{третьей главе} приведены детали реализации платформы, включая структуру проекта, описание API, моделей данных, детерминированного генератора Python-кода, SciPy-симуляции, постпроцессинга уравнений, процесса генерации и результаты тестирования.

В \textbf{четвёртой главе} обсуждаются ограничения текущей реализации (включая ограничения бесплатных языковых моделей), планируемые направления развития и потенциальные области применения платфорформы.

В \textbf{заключении} подводятся итоги работы и формулируются выводы.
